\chapter{Materiales y Métodos}
\label{cap:metodos}

Este capítulo describe el enfoque metodológico del trabajo, siguiendo la metodología de Design Science Research (DSR) de Peffers et al.~\cite{peffers2007}, articulada con los criterios de calidad de Hevner et al.~\cite{hevner2004} y con los estándares empíricos de la comunidad \cite{ralph2021}. La orientación metodológica complementaria se acoge a Runeson y H\"ost para el estudio de caso \cite{runeson2009}, Wohlin et al.~para la experimentación \cite{wohlin2012} y Baltes y Ralph para el muestreo \cite{baltes2022}.

\section{Enfoque metodológico global}
\label{sec:dsr}

El trabajo se instancia en las seis actividades de DSR \cite{peffers2007}:

\begin{enumerate}
  \item \textbf{Identificación del problema:} La cooperativa de transporte carece de un sistema integrado de gestión (Capítulo~\ref{cap:introduccion}, Sección~\ref{sec:problema}).

  \item \textbf{Definición de los objetivos de la solución:} El sistema SGROAS debe gestionar conductores, usuarios, rutas, vehículos, asignaciones e incidentes, con autenticación JWT, 70\% de cobertura mínima y p95 $<$ 200\,ms.

  \item \textbf{Diseño y desarrollo:} Arquitectura de tres capas (Angular + Spring Boot + PostgreSQL + Redis), documentada con ADRs y diagramas C4 (Capítulos~\ref{cap:diseno} y~\ref{cap:implementacion}).

  \item \textbf{Demonstración:} El sistema funcional se despliega con Docker Compose y se verifica con pruebas automatizadas, k6, Lighthouse y OWASP ZAP (Capítulo~\ref{cap:evaluacion}).

  \item \textbf{Evaluación:} Evaluación empírica con métricas de rendimiento, seguridad, usabilidad, cobertura y calidad web (Capítulo~\ref{cap:evaluacion}).

  \item \textbf{Comunicación:} El presente informe y el repositorio público con DOI en Zenodo.
\end{enumerate}

\section{Goal-Question-Metric (GQM)}
\label{sec:gqm}

El diseño empírico se formaliza con el enfoque GQM de Basili et al.~\cite{basili1994}. La meta, preguntas y métricas se resumen en la Tabla~\ref{tab:gqm}.

\begin{longtable}{@{}p{3cm}p{4.5cm}p{5cm}@{}}
\caption{Goal-Question-Metric breakdown of the empirical study.}
\label{tab:gqm}\\
\toprule
\textbf{Meta} & \textbf{Pregunta} & \textbf{Métrica} \\
\midrule
\endfirsthead
\toprule
\textbf{Meta} & \textbf{Pregunta} & \textbf{Métrica} \\
\midrule
\endhead
G1: Rendimiento & RQ1: ¿El acceso híbrido CRUD/ORM + SP mantiene latencia aceptable bajo carga? & p95, media, DT, IC 95\% de \texttt{http\_req\_duration}; tasa de error; throughput \\
\midrule
G2: Seguridad & RQ2: ¿El sistema mitiga los riesgos OWASP Top 10 con controles verificables? & Controles A01--A09 verificados; conteo de SQL dinámico (audit-sql-dynamic.sh) \\
\midrule
G3: Usabilidad & RQ3: ¿Qué nivel de usabilidad perciben los usuarios? & SUS media, DT, IC 95\%; escala adjetiva de Bangor \\
\midrule
G4: Calidad web & RQ4: ¿La interfaz cumple estándares de rendimiento, accesibilidad y SEO? & Lighthouse: Performance, Accessibility, Best Practices, SEO \\
\bottomrule
\end{longtable}

\section{Instrumentos de medición}
\label{sec:instrumentos}

Cada instrumento se describe con la versión exacta y la configuración utilizada:

\begin{itemize}
  \item \textbf{k6 (v0.57):} Pruebas de carga contra \texttt{GET /api/conductores} con 50 VUs y 30\,s. Script versionado en \texttt{k6/script.js} con opciones en \texttt{k6/opts.js}. Datos crudos en \texttt{docs/mediciones/perf/k01--k03.json}.

  \item \textbf{Lighthouse (v12):} Auditoría de calidad web con configuración \texttt{lighthouserc.js}. Perfiles mobile y desktop. Datos en \texttt{docs/mediciones/lighthouse/}.

  \item \textbf{JaCoCo (v0.8.14):} Cobertura de código con reporte XML y HTML generado por \texttt{./mvnw verify}. Datos en \texttt{docs/mediciones/jacoco/jacoco.csv}.

  \item \textbf{OWASP ZAP:} Escaneo baseline de seguridad. Script en \texttt{scripts/zap/run-zap.sh}. Reporte en \texttt{docs/mediciones/sec/zap/}.

  \item \textbf{System Usability Scale (SUS):} Cuestionario de Brooke \cite{brooke1996} con los 10 ítems originales en escala Likert 1--5. Protocolo en \texttt{docs/checklists/ralph2021-<estandar>.md}.

  \item \textbf{SpotBugs/Find-Sec-Bugs:} Análisis estático de seguridad en el pipeline de CI. Regla \texttt{SQL\_NONCONSTANT\_STRING\_PASSED\_TO\_EXECUTE}. Reporte en \texttt{docs/mediciones/sec/static-analysis/}.

  \item \textbf{audit-sql-dynamic.sh:} Script versionado que rechaza SQL dinámico por concatenación en \texttt{db/procs/*.sql} y en el código fuente Java.
\end{itemize}

\section{Población, muestreo y participantes}
\label{sec:muestreo}

La evaluación involucra dos componentes con poblaciones diferentes:

\subsection{Sistema (rendimiento, seguridad, cobertura, calidad web)}

La población de referencia son los endpoints CRUD del sistema SGROAS. Se seleccionó el endpoint \texttt{GET /api/conductores} como representativo por ser el más consultado y el que involucra JOIN con tablas relacionadas. El muestreo es intencional \cite{baltes2022}: se elige un endpoint representativo en lugar de muestrear aleatoriamente entre todos los endpoints, dado el alcance de un PFC.

\subsection{Participantes SUS (usabilidad)}

Participaron 10 personas (P01--P10) anonimizadas, reclutadas por muestreo de conveniencia entre estudiantes y personal de la cooperativa. El perfil declarado incluye: 5 hombres y 5 mujeres, edades entre 19 y 25 años, experiencia web de nivel bajo (1), medio (7) y alto (2), y todos utilizaron computadora de escritorio. El tamaño muestral (n = 10) se ubica en el rango recomendado por Kortum y Bangor \cite{kortum2013}. Los consentimientos informados se archivan fuera del repositorio público, conforme a \texttt{docs/etica/ETHICS.md}.

Originalmente se recolectaron respuestas adicionales de 5 personas más (identificadas como P11--P15), quienes sí dieron consentimiento informado real y verificado. Sin embargo, sus respuestas al cuestionario SUS no contaron con respaldo documental verificable de la fecha en que se declaró que fueron recolectadas, por lo que se excluyeron del análisis oficial para evitar presentar datos sin trazabilidad suficiente (ver \texttt{dataset/sus/PARTICIPANTES-NO-INCLUIDOS.md}).

La Tabla~\ref{tab:sus-demografia} cruza demografía (edad, sexo, experiencia web) con el puntaje SUS por participante; los datos se versionan en \texttt{dataset/sus/sus-raw.csv} y \texttt{P01.json..P10.json}, y los 10 códigos se corresponden 1:1 con el capítulo de resultados.

\begin{table}[htbp]
\caption{SUS demographics and score by participant (P01--P10). Source: \texttt{dataset/sus/sus-raw.csv}.}
\label{tab:sus-demografia}
\small
\begin{tabular}{@{}llllr@{}}
\toprule
\textbf{Código} & \textbf{Edad} & \textbf{Sexo} & \textbf{Exp. web} & \textbf{SUS} \\
\midrule
P01 & 23 & Masculino & Baja & 47.5 \\
P02 & 21 & Femenino & Media & 52.5 \\
P03 & 25 & Masculino & Alta & 50 \\
P04 & 23 & Femenino & Media & 52.5 \\
P05 & 20 & Femenino & Media & 55 \\
P06 & 22 & Masculino & Media & 70 \\
P07 & 20 & Femenino & Media & 90 \\
P08 & 20 & Femenino & Media & 65 \\
P09 & 21 & Masculino & Alta & 77.5 \\
P10 & 19 & Masculino & Media & 70 \\
\bottomrule
\end{tabular}
\end{table}

\section{Protocolo experimental}
\label{sec:protocolo}

\subsection{Rendimiento (k6)}

\begin{enumerate}
  \item Se ejecutan tres corridas independientes de k6 con 50 VUs durante 30\,s contra el endpoint \texttt{GET /api/conductores} protegido con JWT.
  \item Cada corrida se inicia después de reiniciar Redis (cache frío) o manteniendo Redis activo (cache caliente), según la condición.
  \item Se exporta el resumen de cada corrida a \texttt{docs/mediciones/perf/k0N-runN.json}.
  \item Se calcula la media de las medias por corrida, la DT, el IC 95\% con t de Student (gl = $n - 1$), y los percentiles p50, p90, p95, p99.
  \item El umbral declarado a priori es p95 $<$ 200\,ms y tasa de error $<$ 1\,\%.
\end{enumerate}

\subsection{Seguridad (OWASP + ZAP)}

\begin{enumerate}
  \item Se verifican seis controles OWASP Top 10 (A01, A02, A03, A05, A07, A09) con evidencia reproducible (curl o script).
  \item Se ejecuta \texttt{audit-sql-dynamic.sh} sobre \texttt{db/procs/*.sql} y el código fuente Java.
  \item Se ejecuta SpotBugs con Find-Sec-Bugs en el pipeline de CI.
  \item Se ejecuta el escaneo baseline de OWASP ZAP contra la URL pública.
\end{enumerate}

\subsection{Usabilidad (SUS)}

\begin{enumerate}
  \item Cada participante completa las 10 tareas del cuestionario SUS \cite{brooke1996} después de interactuar con el sistema durante 15 minutos.
  \item Las respuestas se registran en escala Likert 1--5.
  \item Se calcula la puntuación SUS: (suma de contribuciones) $\times$ 2.5.
  \item Se reporta media, DT e IC 95\% con t de Student.
  \item Se interpreta con la escala adjetiva de Bangor et al.~\cite{bangor2009,bangor2008}.
\end{enumerate}

\subsection{Cobertura (JaCoCo)}

\begin{enumerate}
  \item Se ejecuta \texttt{./mvnw clean verify} que ejecuta las 293 pruebas JUnit 5 y genera el reporte JaCoCo.
  \item Se extraen los contadores de instrucciones, ramas y líneas del CSV exportado.
  \item El umbral declarado es líneas $\geq$ 70\,\%.
\end{enumerate}

\subsection{Calidad web (Lighthouse)}

\begin{enumerate}
  \item Se ejecutan al menos tres corridas de Lighthouse por perfil (mobile y desktop) contra el frontend Angular.
  \item Se exportan los resultados JSON a \texttt{docs/mediciones/lighthouse/}.
  \item Se reportan las cuatro categorías: Performance, Accessibility, Best Practices, SEO.
  \item Los umbrales declarados son: Performance $\geq$ 80, Accessibility $>$ 90, Best Practices $\geq$ 90, SEO $\geq$ 90.
\end{enumerate}

\section{Análisis estadístico}
\label{sec:estadistica}

Todos los datos cuantitativos presentan media, desviación estándar (DT) e intervalo de confianza al 95\,\% (IC 95\%).

Para el contraste entre condiciones (cache caliente vs cache frío), se utiliza el test de Wilcoxon de muestras pareadas \cite{wilcoxon1945} como test no paramétrico por defecto, acompañado del tamaño de efecto d de Cliff \cite{cliff1993}. Cuando se comparan más de dos condiciones, se aplica la corrección de Holm \cite{wohlin2012}.

El IC 95\% se calcula con la distribución t de Student \cite{student1908}, appropriada para muestras pequeñas ($n < 30$).

\section{Semillas y determinismo}
\label{sec:semillas}

Las versiones exactas de las herramientas se capturan en \texttt{docs/entorno/versions.txt}. Las semillas aleatorias usadas en k6 se declaran en \texttt{k6/opts.js}. Los scripts de análisis (Python) se ejecutan con semillas fijas para garantizar reproducibilidad.

\section{Protocolo de reproducibilidad}
\label{sec:reproducibilidad}

El protocolo completo de reproducibilidad se describe en \texttt{docs/REPRODUCIR.md}. La secuencia end-to-end es:

\begin{verbatim}
git clone https://github.com/Alxjandr07/SGROAS-ProyectoAppWeb.git
cd SGROAS-ProyectoAppWeb
git checkout v1.0.0
make all
\end{verbatim}

El comando \texttt{make all} ejecuta: levantar contenedores (\texttt{docker compose up --build}), pruebas (\texttt{./mvnw test}), benchmarks k6, auditorías estáticas, regeneración de reportes JaCoCo y generación de versiones.
